\section*{Chapter \ref{ch:rs:half-life-decay-and-stability} --- Half-Life, Decay and Stability}

\begin{solution}{exo:rs:half-life-decay-and-stability:1}
$-\ln 2/\ln 0.98 = 34.3$ days; a rank-weighted book trades about $1 - 0.98 = 2\%$ of itself a day.
\end{solution}

\begin{solution}{exo:rs:half-life-decay-and-stability:2}
$1 - 0.37/0.85 = 56\%$ (56.0\% on the unrounded means), close to McLean and Pontiff's average post-publication decline of 58\%.
\end{solution}

\begin{solution}{exo:rs:half-life-decay-and-stability:3}
$0.948(\sqrt{94} + 2/\sqrt{94}) = 9.39$ at the first residual and $0.948(\sqrt{94} + 2 \times 94/\sqrt{94}) = 0.948 \times 3\sqrt{94} = 27.57$
at the last.
\end{solution}

\begin{solution}{exo:rs:half-life-decay-and-stability:4}
A signal of age $a$ still drifts on day $t + h$ if $a + h < 60$, which for $a$ uniform on $0, \dots, 59$ has probability $(60 - h)/60$. The
single-day IC at lag $h$ is proportional to the share of signals still drifting (each drifts at the same rate), so it falls linearly and reaches
zero at $h = 60$. An exponential has no end and a constant halving time; forced onto a straight line that stops, it reports a half-life that
describes neither the start nor the end.
\end{solution}

\begin{solution}{exo:rs:half-life-decay-and-stability:5}
$t = 0.04/(0.10/\sqrt{60}) = 3.10$. By the curve of \cref{fig:rs:half-life-decay-and-stability:posterior}, a measured halving then has a posterior of
about 0.86 (\texttt{rs\_decay.power(0.04, 0.10)} gives 0.858).
\end{solution}

\begin{solution}{exo:rs:half-life-decay-and-stability:6}
Each month the factor's realised return moves all loaded stocks together, in the direction of the \omterm{def:rs:anatomy-of-a-predictor:predictor}{predictor}'s ranks or against them; the rank
correlation with returns then swings with the factor, whatever the \omterm{def:rs:anatomy-of-a-predictor:predictor}{predictor}'s stock-specific information. The swings are common to the whole
cross-section, so they do not average out over stocks, only over months.
\end{solution}

\begin{solution}{exo:rs:half-life-decay-and-stability:7}
\texttt{rs\_decay.detection(0.5)}: sup-F detects a 50\% cut in three markets of ten and the CUSUM in one, between the 70\% cut (seven and three)
and no change at all (one and one, the false alarms). Over a ten-year market with four years after the break, halving a \omterm{def:rs:anatomy-of-a-predictor:predictor}{predictor}'s drift is found
less often than not.
\end{solution}

\begin{solution}{exo:rs:half-life-decay-and-stability:8}
A 24-month IC of momentum is dominated by the factor's realised returns: on the synthetic market the monthly IC has a standard deviation of 0.25
against a mean of 0.02, so two years say almost nothing (a halving is barely evidence at a five-year $t$ of 0.72, and a third over two years even
less). The CUSUM crossing in the last month is exactly what the synthetic momentum did with nothing changed. Before switching off: the test's
false-alarm rate across all monitored \omterm{def:rs:anatomy-of-a-predictor:predictor}{predictors}, the mechanism, crowding and capacity evidence (chapter~28), and a longer sample.
\end{solution}

\begin{solution}{pb:rs:half-life-decay-and-stability:1}
\begin{enumerate}
\item Reversal 0.037 and $-0.04$; surprise 0.010 and 0.983; momentum 0.008 and 0.995; book-to-price 0.004 and 0.9995.
\item Its signal is today's return, uncorrelated with yesterday's ($\rho = -0.04$): the ranks are new every day, and a book that follows them is
rebuilt every day (turnover $1 - \rho = 104\%$).
\item 20 days; the planted drift lasts 60 days.
\item Their single-day ICs (0.004 to 0.008) are near their sampling noise, and a half-life of months cannot show within 60 days.
\item Momentum: mean 0.023, standard deviation 0.247. Surprise: 0.040 and 0.041.
\item From $-0.035$ to 0.071.
\item It loads on the momentum factor, whose monthly return moves the whole cross-section.
\item Both reject at 5\%: the CUSUM crosses in the last month, and sup-F has a permutation p-value of 0.023, with nothing changed.
\item 0.032 before and 0.005 after.
\item At month 53, statistic 12.5, p-value 0.013. Old announcements keep drifting for up to 60 days after the break, so the IC falls over three
months.
\item Nothing: its path reaches 0.69 of the boundary.
\item In this market neither test sees it (sup-F p-value 0.93); over ten markets sup-F detects it in seven and the CUSUM in three.
\item SMB 97\%, HML 55\%, Mom 56\%; sup-F p-values 0.18, 0.15 and 0.053.
\item $p_0 = 0.37$ and $p_1 = 0.50$.
\item \textbf{Named result.} With a prior of one half, a measured halving over five years means true decay with probability 0.57 for a \omterm{def:rs:anatomy-of-a-predictor:predictor}{predictor}
with the synthetic momentum's noise ($t = 0.72$), 0.72 with the momentum factor's pre-publication noise ($t = 1.93$), and 0.999 for one as steady as the
synthetic surprise ($t = 7.6$).
\item About 3.6: with $p_1 \approx \frac12$ the posterior is $0.5/(0.5 + p_0)$, which reaches 0.9 when $p_0 = 0.056$, that is
$\Phi(-t/\sqrt5) = 0.056$ (Interview question~6), $t = 3.6$.
\item The mechanism (is there a reason it should have stopped?), crowding and capacity measures, the \omterm{def:rs:anatomy-of-a-predictor:predictor}{predictor}'s correlation with newly published
ones, trading costs, and the same \omterm{def:rs:anatomy-of-a-predictor:predictor}{predictor}'s behaviour in other markets.
\item Control the false alarms across \omterm{def:rs:anatomy-of-a-predictor:predictor}{predictors} (a stricter level, or a false-discovery rate, chapter~20), require persistence (two consecutive
alarms), and act on size before switching off: reduce the weight when the evidence is weak, remove when it is strong.
\item The original sample's selection (a false or overstated discovery regresses), changes in market structure (decimalisation, liquidity,
costs), and risk that was priced and is priced less. McLean and Pontiff separate the first from trading by comparing out-of-sample and
post-publication periods.
\item When the \omterm{def:rs:anatomy-of-a-predictor:predictor}{predictor} is precise enough that a decline of that size would rarely happen by chance.
\end{enumerate}
\end{solution}

\begin{solution}{iq:rs:half-life-decay-and-stability:1}
The lag at which the signal's autocorrelation falls to one half, $-\ln 2/\ln\rho$ for a lag-one autocorrelation $\rho$. A book that tracks the
signal trades about $1 - \rho$ of itself each period, so short half-lives mean high turnover and costs; the information's own half-life, from the
\omterm{def:rs:half-life-decay-and-stability:decay}{IC decay curve}, sets the holding period, and the two need not agree.
\end{solution}

\begin{solution}{iq:rs:half-life-decay-and-stability:2}
Compute its IC with the return of each single future period $t + h$ (not cumulative returns, which mix horizons), average over many dates, and plot
against $h$; fit a form only after looking at the curve, with bootstrap intervals over dates.
\end{solution}

\begin{solution}{iq:rs:half-life-decay-and-stability:3}
Compare the year with the IC's own noise: its $t$ statistic per year, the probability of such a year with nothing changed, and a break test with
a p-value that accounts for heavy tails (permutation) and for the number of signals monitored. Then look for other evidence: mechanism, crowding,
capacity. A single bad year of a noisy signal is weak evidence.
\end{solution}

\begin{solution}{iq:rs:half-life-decay-and-stability:4}
Statistical bias in the original discovery (it regresses out of sample), investors trading on the publication, and changes in market structure
or costs unrelated to publication. McLean and Pontiff tell the first two apart with the out-of-sample but pre-publication period (a 26\% decline,
an upper bound on data mining) against the post-publication period (58\%); dating the decline against structural changes (decimalisation) tests the
third.
\end{solution}

\begin{solution}{iq:rs:half-life-decay-and-stability:5}
The CUSUM cumulates standardised recursive residuals and rejects when the path leaves boundaries that widen with time; it needs no candidate date
and is cheap, but on this chapter's planted removal it did not cross. Sup-F tries every date, is powerful against a single shift in the mean, and
locates it; its critical value must account for the search (Andrews) or be computed by permutation. Both are one-break tests of the mean and
both reject stable noisy series at their nominal rate.
\end{solution}

\begin{solution}{iq:rs:half-life-decay-and-stability:6}
Let the two means be $\bar x_1 = \mu + \sigma_m z_1$ and $\bar x_2 = \mu + \sigma_m z_2$ with $\sigma_m = \sigma/\sqrt n$ and independent standard
normals. The event $\bar x_2 \le \bar x_1/2$ is $z_2 - z_1/2 \le -\mu/(2\sigma_m)$, and $z_2 - z_1/2$ is normal with variance $5/4$, so the probability is
$\Phi(-t/\sqrt5)$ with $t = \mu/\sigma_m$: 0.37 at $t = 0.72$, 0.19 at $t = 1.93$, 0.0004 at $t = 7.6$.
\end{solution}
